% Shared preamble of the HSEMHS mode documents (docs/physics/*.tex, docs/technical/*.tex).
\documentclass[10pt]{extarticle}
\usepackage[a4paper,margin=1.7cm]{geometry}
\usepackage[T1]{fontenc}\usepackage{lmodern}\usepackage{tabularx}
\usepackage{amsmath,amssymb,bm,enumitem,titlesec,xcolor,booktabs,array,tikz,graphicx,longtable}
\usepackage[hidelinks]{hyperref}
\titleformat{\section}{\large\bfseries\color{blue!50!black}}{\thesection.}{0.5em}{}
\titlespacing*{\section}{0pt}{9pt}{3pt}
\titleformat{\subsection}{\normalsize\bfseries}{\thesubsection}{0.5em}{}
\titlespacing*{\subsection}{0pt}{6pt}{2pt}
\setlist{nosep,leftmargin=1.3em}
\setlength{\parindent}{0pt}\setlength{\parskip}{4pt}
\renewcommand{\arraystretch}{1.15}
\newcommand{\B}{\mathbf B}\newcommand{\J}{\mathbf J}\newcommand{\F}{\mathbf F}\newcommand{\R}{\mathbf R}
\newcommand{\zh}{\hat{\mathbf z}}\newcommand{\Bo}{\mathbf B_{\rm obs}}\newcommand{\Lv}{\mathbf L}
\newcommand{\code}[1]{\texttt{#1}}
\newcolumntype{L}[1]{>{\raggedright\arraybackslash}p{#1}}
\newcolumntype{Y}{>{\raggedright\arraybackslash}X}
\newcommand{\doctitle}[2]{{\Large\bfseries #1}\\[2pt]{\small #2}\par\vspace{2pt}}
\newif\ifadjointcolumn\newif\ifadjointpoisson   % which part of common/adjoint.tex a document shows
\pagestyle{plain}
\AtBeginDocument{\sloppy}

\begin{document}
\doctitle{mhs-poisson-logm: implementation}{Technical notes on \code{run.mode = "mhs-poisson-logm"}. HSEMHS, 25 Nov 2026 (main 9280809).
The solver of \code{mhs-poisson-ltau} (\code{solve\_tau}, \code{mass = true}); read its notes for everything not listed here.}

\section{Running it}
{\footnotesize\begin{verbatim}
[mpiexec -n R] julia --project=. -t T scripts/run.jl config/modes/mhs-poisson-logm.toml
\end{verbatim}}
Driver \code{run\_mhs\_poisson\_ltau} with the mass axis; embedded as \code{prepare\_embedded(comm, :mhs\_poisson\_logm; axis = log10m,
\dots, poisson = PoissonMassOptions(logm\_cut = -0.5, \dots))}. \code{PoissonMassOptions} wraps \code{PoissonTauOptions}, so that a
$\tau$ cut cannot reach a mass axis. Keys: those of \code{mhs-poisson-ltau}, with \code{poisson.logm\_cut} ($-0.5$, $\log_{10}m$ in
kg\,m$^{-2}$) in place of \code{logtau\_cut}; \code{cut\_frame = "vertical"} places it at $\log_{10}m_{\rm LOS}=$ \code{logm\_cut}
$-\log_{10}\mu$. Input, output and units as \code{column-hse-logm} (\code{/mass\_grid}).

\section{What differs from the $\tau$ axis}
\begin{itemize}
\item \code{evaluate!}: $a=\mu m/\rho$ in place of $\mu\tau/\chi$, so the heights $z_{k+1}=z_k-\tfrac12\Delta\ln m_k(a_k+a_{k+1})$
need $\rho$ only; in the Jacobian the opacity derivative $\partial\ln\chi/\partial\ln P$ is replaced by $\partial\ln\rho/\partial\ln P$
(\code{pr.mass}: $\rho$ in the place of $\chi$). Residuals, weights, preconditioners, steps, stop rule and status are unchanged.
\item The start is \code{column\_hse\_tau(\dots; mass = true)}; above the cut $P=\mu gm$ and the start's relative heights.
\item After the solve, $\log\tau$ by \code{logtau\_mass} and the gauge on it. Per-band diagnostics are by $\log m$
(bands $[-9,-1)$, $[-1,0)$, $[0,1)$, $[1,2)$, $[2,9)$).
\end{itemize}

% Shared section of the technical documents of the slab modes (column-hse-ltau, column-hse-logm, mhs-poisson-z,
% mhs-poisson-ltau, mhs-poisson-logm, and the embedded column-hse-z): MPI layer, threads, layout, multi-node runs.
\section{Parallelisation}\label{sec:parallel}
\textbf{MPI layer} (\code{src/Slab.jl}, \code{src/Parallel.jl}).
\begin{itemize}
\item \code{slab\_par(comm, NX, NY, n)}: each rank owns a contiguous range of $x$ columns (\code{slab\_range(NX, R, r)}) and all of
$y$ and the vertical; arrays are (vertical, $y$, $x$), columns contiguous; at least two $x$ columns per rank. The caller's
communicator is duplicated (\code{Comm\_dup}), so HSEMHS messages cannot match the caller's; \code{free\_embedded!} releases it.
MPI runs as \code{THREAD\_FUNNELED}: every MPI call is on the main task, outside threaded loops.
\item $x$ halo planes for every $D_x$ (\code{MPI.Sendrecv!}); global sums, dot products and norms by \code{Allreduce}
(\code{gsum}, \code{gdot}, \code{gnorm}, \code{gmax}).
\item \code{gquantiles}/\code{gmedians}: \emph{exact} global order statistics without gathering. Each Float64 is mapped to an
order-preserving Int64 key; for all sets at once the $k$-th smallest key is found by bisection on the key range, one
\code{Allreduce} of the per-set counts per round (at most 64 rounds). The layer sets are formed in one threaded pass over the
columns and sorted in parallel; the per-round counts stay serial (a threaded loop in each round doubled the allocation of a
column solve).
\item Distributed FFT (\code{SlabFFT}, the Fourier preconditioners): real FFT along $y$, \code{Alltoallv} transpose to slabs of
$k_y$, complex FFT along $x$, and back. Transforms run per (vertical plane, chunk of \code{FFT\_ROWS} = 4 rows), so a rank with
few planes still uses all its threads; the chunking depends only on the grid. FFTW plans are single-threaded, applied from Julia
threads (FFTW's own threading crashed under Julia 1.13).
\item Output: cooperative (rank 0 creates every dataset and attribute, the ranks then write their $x$ hyperslabs in turn), or
parallel HDF5 with collective MPI-IO when \code{output.parallel\_hdf5 = true}, the run has more than one rank, the HDF5
library has MPI-IO and the output directory is on a file system listed in \code{output.parallel\_fs} (Lustre, GPFS, BeeGFS,
\dots; not NFS or StorNext); otherwise the cooperative path. Rank 0 decides, and both paths give the same file.
\end{itemize}

\textbf{Threads} (25 Nov 2026, commit 653a8f5). Every loop over columns is \code{Slab.@cols i j nx ny}: a threaded loop over all
$(j,i)$ columns of the slab, split into 8 chunks per thread, the body a closure behind a function barrier. Earlier the threads only
split the rank's own $x$ columns (252 on one rank, 8 on 32 ranks), so 64--256 threads had little to do. \code{tmap!}, \code{tsum},
\code{tdot} thread element-wise updates and sums in fixed chunks, so a sum has the same value on any thread count. In the Poisson
solvers the operator products (\code{Aop!}, \code{Jop!}, \code{JTop!}), the dense products of the Fourier preconditioner (in column
chunks), the per-mode blocks of the height-grid preconditioner (every $(k_x,k_y)$ a work item) and the pack/unpack of the
transposes are threaded, and the operator and residual buffers are pooled on the solver state instead of allocated per product.
Results and work counters (CG iterations, operator products) are unchanged by the threading. BLAS: one thread by default
(\code{scripts/run.jl}); the standalone $\tau$-grid driver raises it to \code{min(32, T)} around the solve; the embedded API
leaves it to the caller.

\textbf{Layout: one rank per NUMA domain.} A hyades node has 8 NUMA domains (Zen\,2 EPYC 7742: 16 cores each; Zen\,3 EPYC 7543:
8 cores each), each with its own memory controller. A rank's arrays live in one domain (first touch by its main thread), so its
threads share one controller: one rank of 16 threads in one domain is as fast as 64 threads unpinned. Run one rank per domain
with its threads on that domain's cores; \code{test/scripts/numa\_bind.sh} pins rank $r$ of a node to domain
$r \bmod$ \code{NUMA\_PER\_NODE} (\code{numactl --cpunodebind --membind}; local rank from MPICH or OpenMPI). Short
\code{mhs-poisson-ltau} solve at $252^2$ (3 Gauss--Newton steps, no polish), hyades2 (Zen\,2):
\begin{center}\small
\begin{tabular}{lc}
\toprule
layout & solve \\ \midrule
before the threading, 1 rank $\times$ 64 threads & 23.8\,s \\
1 rank $\times$ 16 threads in one domain & 15.9\,s \\
2 / 4 / 8 ranks $\times$ 16, one domain each & 9.8 / 5.8 / 3.8\,s \\
\bottomrule
\end{tabular}
\end{center}
Full solves (polish 3, the adjoint benchmark's settings), one Zen\,3 node, 8 ranks $\times$ 8 threads: \code{mhs-poisson-ltau}
67\,s solve / 52\,s adjoint (was 631 / 555\,s on 1 $\times$ 64), \code{mhs-poisson-logm} 168 / 62\,s, \code{mhs-poisson-z}
12.8 / 2.2\,s (was 281 / 9.5\,s). Most of the $\tau$ and mass-axis time is the polish (3 Newton steps, CG to $10^{-10}$).

\textbf{Across nodes.} MPI.jl's bundled MPICH has no InfiniBand provider and runs over TCP (151\,$\mu$s one-way latency on the
IPoIB interface); across hyades nodes it gave no gain. \code{test/scripts/poisson\_scaling\_ompi.sh} runs over InfiniBand:
OpenMPI 5.0.8 built against the nodes' UCX (\code{--with-ucx=/usr --without-hcoll}, with plain \code{gcc}), a Julia environment
whose MPI.jl uses it (\code{MPIPreferences} system binary; HDF5.jl on the system's serial libhdf5, since HDF5\_jll's OpenMPI build
loads a second libmpi), \code{mpiexec --mca pml ucx}, one rank per domain through \code{numa\_bind.sh}; ping-pong 1.2\,$\mu$s one
way, 24\,GB/s. \code{mhs-poisson-ltau} at $252^2$ (25 Gauss--Newton steps + polish 3), 8 ranks $\times$ 8 threads per Zen\,3 node:
\begin{center}\small
\begin{tabular}{lcccc}
\toprule
nodes (ranks) & solve & adjoint & total & speed-up \\ \midrule
1 (8) & 63.5\,s & 51.0\,s & 114\,s & 1 \\
2 (16) & 35.6\,s & 31.2\,s & 67\,s & 1.7 \\
4 (32) & 28.1\,s & 20.0\,s & 48\,s & 2.4 \\
8 (64) & 23.6\,s & 18.1\,s & 42\,s & 2.7 \\
\bottomrule
\end{tabular}
\end{center}
The work is the same on every layout (76 CG iterations, 6 evaluations). The polish and the adjoint scale; the Gauss--Newton steps
do not (2.6\,s on 1 node, 4.0\,s on 8): a profile of rank 0 on 8 nodes spends their CG almost entirely in MPI, in the
\code{Alltoallv} of the preconditioner's transposes ($\sim$13\,ms each) and the scalar reductions of CG ($\sim$5\,ms each). Whether
that is the collectives at 64 ranks or waiting on slower ranks has not been separated (open). Report:
\code{HSEMHS/reports/poisson\_parallel.md}.

\textbf{Reproducibility.} Medians are exact, but sums are reduced in a rank-dependent order, so results on different rank counts
agree to rounding, and an unconverged iterative solve can amplify that along its path. The column modes have no such
amplification.


\adjointpoissontrue
% Shared section of the technical documents of the six fixed-B modes: the adjoint of the embedded solve.
% \input by each mode document (after \adjointcolumntrue or \adjointpoissontrue, see ../preamble.tex).
\section{Adjoint (embedded API)}\label{sec:adjointapi}
The cotangents of a solve's outputs pulled back to its inputs ($\ln T$ and the cell-centred $\B$), at a cost independent of
how many inputs a caller varies. Design and decisions: \code{design/adjoint/}; results: \code{HSEMHS/reports/adjoint\_phase1-3.md},
\code{adjoint\_phase4-6.md}. Code: \code{src/Adjoint.jl} (API), \code{src/ColumnHSEAdjoint.jl}, \code{src/PoissonTauAdjoint.jl},
\code{src/PoissonZAdjoint.jl}, \code{src/Tau.jl}, \code{src/ViewTau.jl}; the table derivatives are WittBifrostEOS's \code{state\_d}.
{\footnotesize\begin{verbatim}
res = solve_embedded!(ws, lnT, Bx, By, Bz; record = true)  # keeps the tape; any later solve ends it
bar = EquilibriumCotangent(ws)                             # zeros on this rank's slab, (vertical, y, x)
bar.rho .= ...; bar.z .= ...       # any of Pgas, lnP, rho, ne, chi500, tau500, logtau500, z
adj = adjoint_embedded!(ws, bar)   # collective; the tape is kept, so several cotangents work
adj.lnT, adj.Bx, adj.By, adj.Bz    # d(loss)/d ln T and d(loss)/dB on the slab
adj.status                         # :exact, :stationary or :not_stationary; details in adj.info
\end{verbatim}}
\begin{itemize}
\item The call must follow a valid \code{solve\_embedded!} with \code{record = true} on the same workspace; a workspace without
a tape, after a later solve, or whose last solve had a given $n_e$ (\code{lnne}) is refused (\code{InvalidInput}, on every rank).
\code{Pgas} and \code{lnP} cotangents may both be given; they are combined. \code{z\_tau1} is a diagnostic and has no cotangent.
\item Keywords: \code{cg\_tol} ($10^{-10}$; the Poisson adjoint solves), \code{cg\_maxit} (1000), \code{gauss\_newton} (false; true
uses the Gauss--Newton Hessian instead, for comparison), \code{stationary\_tol} ($10^{-6}$).
\ifadjointcolumn
\item \textbf{Column modes:} reverse sweeps of the recorded integrator (\code{RK4Tape} on the height grid, \code{HeunTape} on the
$\tau$ and mass axes: substep counts, the bound that set each Heun step and how often it was quartered, table cells), the top
closure by the implicit function theorem, the alignment with its median reference (\code{gquantile\_vjp}), the gauge, the
$\tau$ integration (\code{tau\_column\_vjp!}, \code{inclined\_tau\_vjp!} with the transposed shear) and the state at the nodes.
\code{adj.status = :exact}. No communication except the reductions of the alignment and the gauge.
\fi
\ifadjointpoisson
\item \textbf{Poisson modes:} the implicit-function derivative at the returned state: one CG solve with the exact Hessian of $\Phi$
(Hessian--vector products by central differences of the exact gradient; relative step $10^{-5}$ on the $\tau$ grid,
\code{HVP\_STEP}), the transposed mixed derivative in $(\ln T,\B)$, the weights through their layer-median columns, the
continuation or fixed part above the cut, the gauge, and the column-HSE start. Directions in which $\Phi$ does not change are
projected out of the solve and their cotangent passed to the start (height grid: the anchor-cell mean of $\ln P$, a bordered
constraint; $\tau$ and mass axes: the uniform height shift and the Nyquist modes of the centred differences). The solve has its
own block-Jacobi preconditioner (the exact Gauss--Newton column blocks with no damping on $\ln P$ and the measured curvature
on $c$); the solver's preconditioners, built for its damped operator, made CG break down. \code{adj.status} is
\code{:stationary} when one undamped Newton step from the returned state predicts a relative decrease of $\Phi$ at most
\code{stationary\_tol} (\code{adj.info["newton\_decrement"]}, solved to CG tolerance $10^{-3}$), else \code{:not\_stationary}.
\item \textbf{\code{polish} = $n$} (\code{PoissonZOptions}, \code{PoissonTauOptions}, default 0; embedding API only, not a
\code{[poisson]} key of the file driver): up to $n$ Newton
steps with the exact Hessian after the Gauss--Newton solve, so that the returned state is the stationary point the adjoint
differentiates. Without it the Gauss--Newton solves can end far from stationarity although their damped decrement is small: on
the $11\times12$ height-grid fixture 80, 300 and 1000 steps gave finite differences 4.4, 7.8 and 10.0 against the adjoint's
10.2, and the anchor mean drifted by $5\times10^{-3}$; 30 steps and 4 Newton steps reach a decrement of $10^{-25}$.
\fi
\end{itemize}

\textbf{Accuracy} (dot-product check against central differences of \code{solve\_embedded!}, Richardson, best over $h$;
\code{test/test\_adjoint.jl}, \code{test/mpi/mpi\_adjoint.jl} on 2 and 3 ranks, which also equal rank 0 alone to $10^{-10}$):
\begin{center}\small
\begin{tabular}{lcc}
\toprule
case & 1 rank & 3 ranks \\ \midrule
\ifadjointcolumn
column-hse-ltau (ptot, $\B$); align none, $\B=0$; $\mu$ 0.8 inclined & $3.8\times10^{-10}$; $1.8\times10^{-9}$; $9.4\times10^{-9}$ & the same \\
column-hse-logm (ptot, $\B$); $\mu$ 0.8 inclined & $9.3\times10^{-10}$; $2.6\times10^{-10}$ & $9.4\times10^{-10}$; $2.5\times10^{-10}$ \\
column-hse-z anchor bottom / top / height; $\mu$ 0.8 inclined & $3.0\times10^{-8}$ / $5.7\times10^{-10}$ / $4.9\times10^{-9}$; $2.3\times10^{-9}$ & the same \\
\fi
\ifadjointpoisson
mhs-poisson-z (30 Gauss--Newton steps + polish) & $6.4\times10^{-8}$ & $3.5\times10^{-8}$ \\
mhs-poisson-ltau, default solver + polish & $3.0\times10^{-8}$ & $1.4\times10^{-9}$ \\
mhs-poisson-ltau $\mu$ 0.8 inclined; block preconditioner & $5.8\times10^{-9}$; $2.3\times10^{-9}$ & $5.8\times10^{-9}$; $2.0\times10^{-9}$ \\
mhs-poisson-logm, default solver + polish & $3.5\times10^{-9}$ & $3.9\times10^{-9}$ \\
\fi
\bottomrule
\end{tabular}
\end{center}
\textbf{Cost} (252$^2$ columns, one rank $\times$ 64 threads, hyades5, \code{test/scripts/adjoint\_benchmark.jl}; the column
times vary by about 30\,\% between runs):
\begin{center}\small
\begin{tabular}{lccc}
\toprule
mode & solve & recorded solve & adjoint \\ \midrule
\ifadjointcolumn
column-hse-ltau (96 levels) & 0.37\,s & 1.24\,s & 0.60\,s \\
column-hse-logm (96 levels) & 0.36\,s & 0.79\,s & 0.63\,s \\
column-hse-z (177 levels) & 0.81\,s & 1.74\,s & 2.51\,s \\
\fi
\ifadjointpoisson
mhs-poisson-ltau (95 levels, polish 3), 1 rank $\times$ 64 threads & 631\,s & 649\,s & 555\,s \\
the same, 8 ranks $\times$ 8 threads on one Zen\,3 node (Sect.~\ref{sec:parallel}) & 67\,s & -- & 52\,s \\
mhs-poisson-z (177 levels, polish 3), 1 $\times$ 64 / 8 $\times$ 8 & 281\,s / 12.8\,s & 275\,s / -- & 9.5\,s / 2.2\,s \\
\fi
\bottomrule
\end{tabular}
\end{center}
Limits: the derivatives at the table's slope breaks (node pairs $10^{-9}$ apart in $\ln T$) are those of the cell the point
falls in, so a finite difference straddling a break does not match; the median weights are exact away from ties.


\section{Verification and results}
\begin{itemize}
\item \code{test\_logm.jl} (``Jacobian on the mass axis'' at $(\mu,\phi)$ = (1, 0), (0.6, 45$^\circ$), (0.25, 120$^\circ$); ``Poisson
solve''; viewing angle), \code{mpi/mpi\_view.jl}, \code{test\_adjoint.jl}, \code{mpi/mpi\_adjoint.jl}.
\item $504^2\times76$ cube, one rank $\times$ 64 threads (\code{HSEMHS/reports/logm\_cut\_2026-11-03}): at the default cut 5 steps,
converged, 167\,s; $z(\tau{=}1)$ rms 19.3\,km, correlation 0.778, Wilson depression $-82.5$\,km (snapshot $-77.9$\,km), rms
$\Delta\ln P$ 0.069. Cuts $-2\ldots+0.5$ agree within 0.3\,km in $z(\tau{=}1)$ rms; $-3$ ends \code{max\_iter} at 25 steps.
\item $252^2$ (\code{IO/tau252/input\_logm\_mu1.0.nc}, \code{paper\_252\_2026-11-10}): 6 steps, converged; $z(\tau{=}1)$ rms
$56.6\to19.5$\,km. With polish 3 on one Zen\,3 node, 8 ranks $\times$ 8 threads: 168\,s solve, 62\,s adjoint.
\item The rows are for an input whose $m$ comes from the snapshot's own $\rho$ and $P$; an inversion's $m$ is not that exact.
\end{itemize}
\end{document}
